% Part:sets-functions-relations
% Chapter: size-of-sets
% Section: reduction-alt

\documentclass[../../../include/open-logic-section]{subfiles}

\begin{document}

\olfileid[fa-IR]{sfr}{siz}{red-alt}

\olsection{کاهش}

\begin{editorial}
  این بخش ناشمارابودن را از راه کاهش ثابت می‌کند و با نتایجِ
  \olref[nen-alt]{sec} مطابقت دارد. روایتی جایگزین و اندکی مشروح‌تر که
  با نتایجِ \olref[nen]{sec} مطابقت دارد، در \olref[red]{sec} آمده است.
\end{editorial}

با استدلالی قطری ثابت کردیم که $\Bin^\omega$ !!{nonenumerable} است.
با استدلال قطریِ مشابهی نشان دادیم که $\Pow{\Nat}$ !!{nonenumerable}
است. اما راه دیگری نیز برای اثبات اینکه $\Pow{\Nat}$
!!{nonenumerable} است وجود دارد: نشان دهید که \emph{اگر
$\Pow{\Nat}$ !!{enumerable} باشد، آنگاه $\Bin^\omega$ نیز
!!{enumerable} است}. چون می‌دانیم $\Bin^\omega$ !!{nonenumerable} است،
نتیجه می‌شود که $\Pow{\Nat}$ نیز چنین است.

این کار را \emph{کاهش‌دادن} یک مسئله به مسئله‌ای دیگر می‌نامند. در این
مورد، مسئلهٔ برشمردن $\Bin^\omega$ را به مسئلهٔ برشمردن
$\Pow{\Nat}$ کاهش می‌دهیم. راه‌حلی برای مسئلهٔ دوم---یعنی برشماریِ
$\Pow{\Nat}$---راه‌حلی برای مسئلهٔ نخست---یعنی برشماریِ
$\Bin^\omega$---به دست می‌دهد.

برای کاهش مسئلهٔ برشمردن مجموعهٔ~$B$ به مسئلهٔ برشمردن مجموعهٔ~$A$،
روشی ارائه می‌کنیم که برشماریِ~$A$ را به برشماریِ~$B$ تبدیل کند.
آسان‌ترین راه برای این کار تعریف یک !!a{surjection} مانند
$f\colon A \to B$ است. اگر $x_1$، $x_2$، \dots{} مجموعهٔ~$A$ را
برشمارد، آنگاه $f(x_1)$، $f(x_2)$، \dots{} مجموعهٔ~$B$ را برمی‌شمارد.
در مورد ما، در پی یک !!a{surjection} مانند $f\colon \Pow{\Nat}
\to \Bin^\omega$ هستیم.

\begin{prob}
نشان دهید اگر یک تابعِ !!{injective} مانند $g\colon B \to A$ وجود داشته
باشد و $B$~!!{nonenumerable} باشد، آنگاه $A$ نیز چنین است. برای این کار
نشان دهید چگونه می‌توان با استفاده از~$g$، برشماریِ~$A$ را به برشماریِ~$B$
تبدیل کرد.
\end{prob}

\begin{proof}[برهانِ {\olref[nen-alt]{thm:nonenum-pownat}} از راه کاهش]
برای انجام کاهش، فرض کنید $\Pow{\Nat}$ !!{enumerable} است و بنابراین
برشماری‌ای از آن به‌صورت $N_{1}$، $N_{2}$، $N_{3}$، \dots وجود دارد.

تابع $f \colon \Pow{\Nat} \to \Bin^\omega$ را با این قرارداد تعریف
کنید که $f(N)$ رشتهٔ $s_{k}$ باشد، به‌گونه‌ای که $s_{k}(n) = 1$ اگر و تنها
اگر $n \in N$، و در غیر این صورت $s_k(n) = 0$.

این امر به‌روشنی تابعی را تعریف می‌کند، زیرا هرگاه $N \subseteq \Nat$،
هر $n \in \Nat$ یا یک !!a{element} از $N$ است یا نیست. برای نمونه،
مجموعهٔ $2\Nat = \Setabs{2n}{n \in \Nat} = \{0,2, 4, 6,
\dots\}$ از اعداد طبیعی زوج به رشتهٔ $1010101\dots$ نگاشته می‌شود؛
$\emptyset$ به $0000\dots$ نگاشته می‌شود و $\Nat$ به
$1111\dots$.

این نگاشت همچنین !!{surjective} است: هر رشته از $0$ها و $1$ها با
مجموعه‌ای از اعداد طبیعی متناظر است؛ یعنی مجموعه‌ای که اعضایش آن اعداد
طبیعی‌اند که با جایگاه‌هایی متناظرند که رشته در آن‌ها $1$ دارد. دقیق‌تر
بگوییم، اگر $s \in \Bin^\omega$، آنگاه $N
\subseteq \Nat$ را چنین تعریف کنید:
\[
N = \Setabs{n \in \Nat}{s(n) = 1}
\]
آنگاه $f(N) = s$؛ این را می‌توان با مراجعه به تعریفِ~$f$ بررسی کرد.

اکنون فهرست زیر را در نظر بگیرید:
\[
f(N_1), f(N_2), f(N_3), \dots
\]
چون $f$ !!{surjective} است، هر عضوِ $\Bin^\omega$ باید برای ورودی‌ای
به‌عنوان مقدارِ~$f$ ظاهر شود و بنابراین باید در فهرست بیاید. پس این
فهرست باید همهٔ~$\Bin^\omega$ را برشمارد.

بنابراین اگر $\Pow{\Nat}$ !!{enumerable} بود، $\Bin^\omega$ نیز
!!{enumerable} می‌بود. اما $\Bin^\omega$ !!{nonenumerable} است
(\olref[nen-alt]{thm:nonenum-bin-omega}). پس $\Pow{\Nat}$
!!{nonenumerable} است.
\end{proof}

%\begin{explain}
%It is easy to be confused about the direction the reduction goes in.
%For instance, !!a{surjective} function $g \colon \Bin^\omega \to X$
%does \emph{not} establish that $X$ is !!{nonenumerable}.  (Consider $g
%\colon \Bin^\omega \to \Bin$ defined by $g(s) = s(1)$, the function
%that maps a sequence of $0$'s and $1$'s to its first !!{element}.  It
%is surjective, because some sequences start with $0$ and some start
%with $1$. But $\Bin$ is finite.)  Note also that the function $f$ must
%be surjective, or otherwise the argument does not go through:
%$f(x_1)$, $f(x_2)$, \dots{} would then not be guaranteed to include
%all the !!{element}s of~$Y$. For instance, $h\colon \Nat \to
%\Bin^\omega$ defined by
%\[
%h(n) = \underbrace{000\dots0}_{\text{$n$ $0$'s}}
%\]
%is a function, but $\Nat$ is !!{enumerable}.
%\end{explain}

\begin{prob}\label{sfr:siz:red-alt:prob:nat-nat}
  با استدلالی کاهشی نشان دهید مجموعهٔ~$X$ از همهٔ توابع
  $f\colon \Nat \to \Nat$ !!{nonenumerable} است. (راهنمایی: یک نگاشت
  پوشا از $X$ به~$\Bin^\omega$ بدهید.)
\end{prob}

\begin{prob}
با استدلالی کاهشی نشان دهید که مجموعهٔ همهٔ \emph{مجموعه‌های} زوج‌های
اعداد طبیعی، یعنی $\Pow{\Nat \times \Nat}$، !!{nonenumerable} است.
\end{prob}

\begin{prob}
با استدلالی کاهشی نشان دهید که $\Nat^\omega$، یعنی مجموعهٔ دنباله‌های
نامتناهیِ اعداد طبیعی، !!{nonenumerable} است.
\end{prob}

%\begin{prob}
%Let $P$ be the set of functions from $\Nat$ to the set $\{0\}$, and let $Q$ be the set of \emph{partial}
%functions from the set of positive integers to the set $\{0\}$. Show
%that $P$~is !!{enumerable} and $Q$~is not. (Hint: reduce the problem
%of enumerating $\Bin^\omega$ to enumerating~$Q$).
%\end{prob}

\begin{prob}
فرض کنید $S$ مجموعهٔ همهٔ !!{surjection}s از $\Nat$ به مجموعهٔ
$\{0,1\}$ باشد؛ یعنی $S$ از همهٔ !!{surjection}s~$f \colon \Nat
\to \Bin$ تشکیل می‌شود. نشان دهید $S$ !!{nonenumerable} است.
\end{prob}

\begin{prob}
نشان دهید مجموعهٔ~$\Real$ از همهٔ اعداد حقیقی !!{nonenumerable} است.
\end{prob}

\end{document}
